% ============================================================
% CHƯƠNG 1: TỔNG QUAN ĐỀ TÀI
% ============================================================
\section{Tổng quan đề tài}

\subsection{Vấn đề đặt ra}

\subsubsection{Thực trạng trì hoãn chuyến bay tại Mỹ năm 2015}

Theo dữ liệu của Cục Hàng không Liên bang Mỹ (FAA), năm 2015 ghi nhận hơn 5,8 triệu chuyến bay nội địa,
trong đó xấp xỉ 21\% bị trì hoãn (delay) với thời gian trễ trên 15 phút.
Tình trạng này gây thiệt hại kinh tế ước tính hàng tỷ đô la mỗi năm, bao gồm chi phí nhiên liệu tăng thêm,
bồi thường hành khách, và ảnh hưởng đến lịch trình của toàn bộ hệ thống hàng không.

Các nguyên nhân chính gây delay được ghi nhận bao gồm:
\begin{itemize}
    \item \textbf{Thời tiết} (Weather Delay): Bão, sương mù, tuyết rơi ảnh hưởng đến tầm nhìn và khả năng cất hạ cánh.
    \item \textbf{Hệ thống quản lý không lưu} (NAS Delay): Tắc nghẽn không phận, hạn chế đường bay.
    \item \textbf{Vận hành của hãng hàng không} (Carrier Delay): Sự cố kỹ thuật, thiếu phi hành đoàn.
    \item \textbf{Tàu bay đến muộn} (Late Aircraft Delay): Hệ quả dây chuyền từ chuyến bay trước bị delay.
    \item \textbf{An ninh} (Security Delay): Kiểm tra an ninh kéo dài.
\end{itemize}

\subsubsection{Hạn chế của các phân tích thủ công truyền thống}

Các phương pháp phân tích truyền thống dựa trên báo cáo tĩnh, bảng tính Excel hoặc truy vấn SQL đơn giản
gặp phải nhiều hạn chế khi xử lý tập dữ liệu ở quy mô hàng triệu bản ghi:
\begin{itemize}
    \item Không thể phân tích đa chiều (theo hãng, sân bay, tháng, ngày trong tuần) một cách linh hoạt.
    \item Thiếu khả năng phát hiện mẫu ẩn (hidden patterns) trong dữ liệu lớn.
    \item Không dự báo được xác suất delay cho chuyến bay mới trước khi xảy ra.
    \item Chi phí xử lý cao và thời gian phân tích kéo dài khi dữ liệu tăng lên.
\end{itemize}

\subsubsection{Nhu cầu có hệ thống phân tích tự động và dự báo}

Để giải quyết bài toán trên, cần một hệ thống tích hợp có khả năng:
\begin{enumerate}
    \item Tự động thu thập, làm sạch và lưu trữ dữ liệu hàng không theo chuẩn Data Warehouse.
    \item Cung cấp giao diện phân tích OLAP đa chiều, cho phép drill-down/roll-up theo nhiều chiều dữ liệu.
    \item Áp dụng các thuật toán Machine Learning để dự báo khả năng delay của chuyến bay mới.
\end{enumerate}

% ------------------------------------------------------------
\subsection{Mục tiêu}

Đề tài hướng đến ba mục tiêu chính:

\begin{enumerate}
    \item \textbf{Xây dựng Data Warehouse chuẩn hóa dữ liệu hàng không:}
    Thiết kế và triển khai kho dữ liệu theo mô hình Star Schema với bảng sự kiện \texttt{fact\_flights}
    và các bảng chiều \texttt{dim\_airline}, \texttt{dim\_airport}, \texttt{dim\_date}, \texttt{dim\_time},
    đảm bảo tính nhất quán và khả năng truy vấn hiệu quả.

    \item \textbf{Cung cấp dashboard phân tích OLAP đa chiều:}
    Xây dựng dashboard tương tác 2 cấp phân tích --- cấp tổng quan (tỷ lệ delay theo hãng, sân bay, tháng)
    và cấp nguyên nhân (phân tích chi tiết từng loại delay) --- cho phép người dùng lọc động theo nhiều tiêu chí.

    \item \textbf{Dự báo khả năng delay chuyến bay mới:}
    So sánh và đánh giá 5 thuật toán phân loại (Logistic Regression, Random Forest, Neural Network,
    XGBoost, LightGBM) trên cùng tập dữ liệu, chọn ra mô hình tốt nhất để dự báo delay.
\end{enumerate}

% ------------------------------------------------------------
\subsection{Phạm vi}

\paragraph{Dữ liệu:}
Đề tài sử dụng bộ dữ liệu \textit{US Flight Delays 2015} công khai trên Kaggle \cite{kaggle2015},
bao gồm ba file CSV:
\begin{itemize}
    \item \texttt{flights.csv}: 5.819.079 chuyến bay, 31 thuộc tính (thời gian, hãng bay, sân bay, các loại delay).
    \item \texttt{airlines.csv}: Thông tin 14 hãng hàng không nội địa Hoa Kỳ.
    \item \texttt{airports.csv}: Thông tin 322 sân bay (tên, thành phố, tiểu bang, tọa độ địa lý).
\end{itemize}

\paragraph{Thời gian và không gian:}
Dữ liệu bao phủ toàn bộ năm 2015 (1/1/2015 -- 31/12/2015), chỉ bao gồm các chuyến bay nội địa Hoa Kỳ.
Dữ liệu quốc tế và các năm khác nằm ngoài phạm vi của đề tài.

% ------------------------------------------------------------
\subsection{Đóng góp}

Đề tài đạt được ba đóng góp kỹ thuật chính:

\begin{enumerate}
    \item \textbf{ETL pipeline tối ưu bộ nhớ:}
    Pipeline xử lý dữ liệu theo kiểu streaming (chunk 100.000 dòng/lần) kết hợp tối ưu dtype
    (float32/int8/int16 thay cho float64/int64) và vectorized transform với \texttt{numpy.select()},
    cho phép nạp toàn bộ 5,8 triệu dòng vào SQLite với mức tiêu thụ RAM chỉ \textasciitilde{}120 MB/chunk,
    thời gian ước tính 3--8 phút tùy tốc độ ổ đĩa.

    \item \textbf{So sánh hệ thống 5 thuật toán Machine Learning:}
    Xây dựng framework đánh giá thống nhất cho 5 thuật toán (Logistic Regression, Random Forest,
    Neural Network, XGBoost, LightGBM) trên cùng tập train/test với các chỉ số AUC-ROC, F1,
    Precision, Recall và thời gian huấn luyện, cho phép so sánh khách quan và toàn diện.

    \item \textbf{Dashboard tương tác 2 cấp phân tích:}
    Giao diện Streamlit 2 tab: (1) tổng quan với 4 KPI cards và biểu đồ theo hãng/sân bay/tháng/quý,
    (2) phân tích nguyên nhân delay với pie chart, delay theo mùa và theo khung giờ bay.
    Bộ lọc động 4 chiều: thời gian, hãng bay, sân bay đi (lọc theo bang), sân bay đến (lọc theo bang).
\end{enumerate}
